\section{3. Background}

\subsection{3.1 Attention as an Integral Operator}

For continuous query and key functions $Q, K: \Omega \to \mathbb{R}^{d_h}$:

\[
(\mathcal{K}Q)(s) = \int_\Omega \kappa(s, t) V(t) \, dt
\tag{3.1}
\]

with $\kappa(s, t) = \exp(\langle Q(s), K(t) \rangle / \sqrt{d_h})$.

Discretizing on a grid $\{t_1, \dots, t_N\}$ and normalizing:

\[
(\mathcal{K}Q)_i \approx \frac{\sum_j \kappa(s_i, t_j) V(t_j)}{\sum_j \kappa(s_i, t_j)}
\tag{3.2}
\]

which is the softmax attention formula.

\subsection{3.2 Galerkin Discretization}

Choose a finite-dimensional subspace $\mathcal{V}_r = \text{span}\{\phi_1, \dots, \phi_r\}$ 
with $\phi_k \in L^2(\Omega)$. The Galerkin approximation $\hat{Q} \in \mathcal{V}_r$ 
satisfies:

\[
\langle \mathcal{K} \hat{Q}, \phi_j \rangle = \langle \mathcal{K} Q, \phi_j \rangle \quad \forall j
\tag{3.3}
\]

For orthonormal $\{\phi_k\}$, the projected operator is:

\[
(\mathcal{K}_r)_{ij} = \langle \mathcal{K} \phi_j, \phi_i \rangle
\tag{3.4}
\]

In matrix form, if $\Phi = [\phi_1, \dots, \phi_r] \in \mathbb{R}^{N \times r}$, 
then $\mathcal{K}_r = \Phi^\top \mathcal{K} \Phi$.

\subsection{3.3 Discrete Cosine Transform (DCT)}

The DCT-II basis:

\[
\phi_k(n) = \sqrt{\frac{2 - \delta_{k0}}{N}} \cos\!\left(\frac{\pi k (2n + 1)}{2N}\right)
\tag{3.5}
\]

\textbf{Properties:}

1. \textbf{Orthonormality:} $\langle \phi_i, \phi_j \rangle = \delta_{ij}$.
2. \textbf{Real-valued:} Unlike FFT.
3. \textbf{Energy compaction:} DCT concentrates most signal energy in the low-frequency coefficients. This is critical for attention, where smooth variations dominate.
4. \textbf{Fast transform:} $O(N \log N)$ via FFT.

\subsection{3.4 Why DCT?}

Attention outputs tend to be \textbf{smooth} over the sequence:
\begin{itemize}
\item The kernel $\kappa(s, t)$ is smooth in both arguments.
\end{itemize}
\begin{itemize}
\item Values $V(t)$ vary slowly.
\end{itemize}
\begin{itemize}
\item The attention output $(\mathcal{K}Q)(s)$ is a smoothed integral.
\end{itemize}

Therefore, a \textbf{low-rank basis that captures smoothness} is ideal. DCT is 
the optimal basis for this under certain smoothness assumptions (Karhunen-Lo ve 
theorem for first-order Markov processes).

\subsection{3.5 Energy Preservation}

When reconstructing from a rank-$r$ projection, the total energy scales as:

\[
\|\Phi \Phi^\top x\|_2^2 = \|\Phi^\top x\|_2^2 \cdot \text{(correction)}
\tag{3.6}
\]

For DCT with $r < N$, a naive reconstruction underestimates the signal 
energy by a factor. We apply an empirical scaling:

\[
\text{Output} = \Phi \cdot O_\Phi \cdot \sqrt{N / r}
\tag{3.7}
\]

to compensate.

\subsection{3.6 Grouped Query Attention}

GQA uses $H_{kv} < H$ KV heads. Galerkin preserves GQA structure via \texttt{repeat\\_kv}.

\subsection{3.7 Rotary Position Embeddings}

RoPE rotates Q, K before the projection step. Partial RoPE is supported via 
coordinate slicing.

\subsection{3.8 Notation Summary}
